\pdfminorversion=7
\documentclass[draft,linenumbers]{agujournal2025}
\usepackage{booktabs,amsmath,amssymb,microtype,url,graphicx,array}
\usepackage{tabularx}
\usepackage{ragged2e}
\setlength{\emergencystretch}{3em}
\input{review_completion_claims.tex}
\input{simple_research_claims.tex}
\input{core_claims_values.tex}
\input{summer_window_values.tex}
\input{beaufort_proxy_values.tex}
\input{summer_proxy_archive_values.tex}
\input{past_season_values.tex}
\input{uncertainty_paired_values.tex}
\input{support_projection_values.tex}
\input{peak_visibility_values.tex}
\input{annual_interval_expansion_values.tex}
\input{nova_scotia_exploratory_values.tex}
\input{nova_clock_sensitivity_values.tex}
\input{nova_field_values.tex}
\input{truckee_values.tex}
\input{region_holdout_values.tex}
\input{au_prior_exposure_values.tex}
\input{source_rights_values.tex}
\input{open_source_values.tex}
\input{open_source_annual_values.tex}
\input{single_year_calibration_values.tex}
\input{single_year_trend_values.tex}
\input{strict_region_values.tex}
\input{input_precision_values.tex}
\input{release_matched_values.tex}
\input{release_regional_values.tex}
\input{archive_matched_values.tex}
\input{expanded_stress_values.tex}
\input{heat_gap_matched_values.tex}
\input{heat_gap_source_values.tex}
\input{stress_source_values.tex}
\input{projected_risk_values.tex}
\input{groupmax_values.tex}
\input{bc_exploratory_values.tex}
\input{eccc_red_values.tex}
\input{external_qualification_values.tex}
\begin{document}
\journalname{Water Resources Research}
\title{Incomplete river-temperature years: observed heat, finite-record retention and the limits of trend inference}
\authors{[Author names]\affil{1}}
\affiliation{1}{[Affiliations]}
\authoraddr{[Corresponding author and email]}
\authorrunninghead{Anonymous}\titlerunninghead{Incomplete river-temperature years}
\keypoints
{Observed windows in incomplete years can overturn a complete-year-only heat record, but most such witnesses also pass a warm-season rule}
{Retaining years improves finite-record slopes; the two calibrated estimators are not uniformly ranked}
{Annual, task and three-output calibration do not certify trend intervals or group risk; strict external validation remains open}
\maketitle\par
\begin{abstract}
Incomplete river-temperature years can retain sustained-heat information, but an annual value, a released number and a trend need different evidence. At \RRecordStations{} of \RRecordDenominator{} long-record stations, fully observed 30-day windows in \RRecordStationYears{} incomplete years exceed every complete-year heat maximum. Of these witnesses, \CCWarmWitnessAdmitted{} pass a station-specific 90\% warm-season rule and \CCAnnualWitnessAdmitted{} pass 95\% annual coverage, so the observation-only result mainly challenges strict completeness. Across twenty source repetitions, annual distribution regression and calibrated stochastic reconstruction have similar errors on difficult gaps. On \CCCohortCommon{} common forward controlled-gap rows, retaining supported years gives a root-mean-square error of \CCAllCommonRMSE{} $^{\circ}$C per decade, compared with \CCWarmRMSE{} under the warm-season rule. Retaining years can improve that finite-record slope, but estimator rank changes with setting and region. Where the reference peak exceeds the observed bound, nominal 90\% coverage is \CCPeakHiddenDTwoCov\% and \CCPeakHiddenStochCov\%. At 4,096 joint-rank samples, nominal 95\% trend coverage remains 87.4--89.9\%. Task calibration is available for \CCTaskPercent\% of forward cases and widens intervals. A three-output rule does not beat score ranking. External checks leave reference quality, timing or watershed dependence unresolved. The results separate observed windows, annual estimates and finite-record trends.
\begin{plainlanguagesummary}
Gaps in a river's water-temperature record do not necessarily erase information about a hot year. If thirty consecutive observed days are warmer than the complete-year record, the complete-year-only record is already incomplete. Most such examples in our archive also pass a less restrictive warm-season rule. We compare ways to estimate missing values and ask what they imply for a trend across many years. A calibrated classical reconstruction is about as accurate as a more elaborate annual model. Trend intervals become less reliable when applied to later years, and a warning that a value may be unsafe does not provide a general group-risk guarantee. The analysis separates observed facts, estimated values and cases where the evidence is insufficient.
\end{plainlanguagesummary}
\end{abstract}

\section{Introduction}
Stream-temperature records support studies of thermal habitat, climate change and water management \cite{caissie2006,michel2020,vanhamel2024}. The annual maximum of a 30-day mean summarizes sustained heat and is related to the thermal peak used in river-network studies \cite{beaufort2022}. A monitoring gap can affect this statistic very differently depending on its timing, the observations remaining in the year and the information available from other years or stations.

Prior work has already shown that partial stream-temperature records can support some annual analyses. \citeA{johnson2021} evaluated missing intervals and imputation for annual mean, phase and amplitude, finding limitations that varied among rivers. \citeA{beaufort2022} used summer observations and climate correction for a multiyear thermal-peak metric. Neither should be interpreted as requiring every calendar year to be complete, and neither supplies a universal error baseline for annual H30. Our question is narrower: what do incomplete daily records support about an annual sustained-heat maximum, numerical publication, and a finite-record trend? We compare strict completeness with a station-specific warm-season rule and distinguish a mathematical observation bound from model-based estimation.

Regression record augmentation and maintenance-of-variance methods have a long hydrological history \cite{matalas1964,hirsch1982,vogel1985}. A point prediction, a predictive distribution and a decision to release a value have different targets. Distributional imputation should be assessed for the downstream quantity rather than only for daily squared error \cite{naf2023,gneiting2007}. Here we compare a daily plug-in, simple corrections, calibrated stochastic reconstruction and annual distribution regression. We then jointly hide portions of a station's multiyear record before local refitting and evaluate whether retaining years changes the slope of that finite record. Finally, we examine whether the annual distributions support trend uncertainty and whether visible information can identify cases that should not receive a numerical estimate.

The hypotheses are deliberately conditional: target-oriented recovery can improve on a daily plug-in without necessarily beating a calibrated stochastic comparator; retaining years can improve a finite-record slope under a specified degradation mechanism; and annual and trend uncertainty require separate evaluation. Table~\ref{tab:core-claims} states the three claims, the common denominator used here and the boundary that limits each claim. The evidence map in the Supporting Information retains the more detailed comparison list.
\input{core_claims_table.tex}

\section{Data and study design}
The source archive contains previously assembled European and US daily observations: 581 physical stations, 185 connected-record groups, 27 regions and 6,170 unique complete station-years. NASA POWER daily air temperature is the common meteorological input \cite{nasapower2026}. Artificial controlled gaps of 30, 60 or 90 days are crossed with overlap with a training-period climatological peak window; empirical incomplete calendars provide a separate family. Each of twenty source repetitions redraws masks and group-role composition. In the leave-one-year-out source comparison, fitting, interval calibration, score learning and rule calibration use disjoint connected groups outside each held region. These repetitions describe variability within the available archive, not twenty independent climates. Figure~\ref{fig:architecture} separates the observation-only, single-year and whole-archive evidence.

We use complete daily station-years as reference outcomes for artificial masking. Water observations hidden by a mask never enter the local fit for that case. Features available when a value is evaluated include the visible gap geometry, a training-period climatological peak, local water history, air conditions and available neighboring stations. The true peak location and hidden water values are evaluation quantities. Observation quality is inherited from the assembled source records; a reported value is not assumed to be error-free physical truth. An observed lower bound is valid for the reported series if its thirty measurements are comparable and valid.

The annual target is
\begin{equation}
 H30_y=\max_{j}\frac{1}{30}\sum_{d=j}^{j+29}Y_{yd}.
\end{equation}
For an incomplete year, the largest fully observed 30-day mean $L_y$ satisfies $H30_y\ge L_y$ for the same reported series. For 1981--2025 we select stations with at least fifteen complete reference years and compare every incomplete-year bound with the largest complete-year H30. The observation-only denominator contains \CCEligibleYears{} eligible station-years at \CCEligibleStations{} stations in \CCEligibleGroups{} connected groups. We evaluate whether the same witnesses pass a 95\% annual-coverage rule and a 90\% station-specific warm-season rule. The five consecutive warm months are obtained from the station's climatology after excluding the evaluated year. These are explicit study comparators; the seven-week missing-period result of \citeA{johnson2021} concerns annual signal parameters, not a validated H30 record rule.

The whole-archive benchmark contains 42 stations with at least twenty complete reference years in seven regions, 18 connected groups and 3,320 station--repetition--setting evaluations. One Finnish record ends in 1994, leaving 41 stations for forward evaluation. For the main retention contrast in the forward controlled-gap setting, both retention sets are available in \CCCohortCommon{} of \CCCohortAll{} settings; the other \CCCohortAll{} minus \CCCohortCommon{} settings are used only for availability accounting. The held-test-group main analysis excludes evaluated test groups from the named source roles. The forward HUC17 estimator and NVE interval calibration retain additional same-region training groups, so the strict-region sensitivity is reported separately. The saved historical \texttt{design.json} files are not rewritten; a correction notice accompanies them.

\input{architecture_figures.tex}

\section{Methods}
\subsection{Annual recovery and the decision to release}
The daily plug-in computes H30 after a selected daily reconstruction. Functional debiasing corrects its annual value, whereas D2 learns annual-error quantiles using visible features and simulated trajectories. A stochastic regression comparator generates conditional residual bridges and calculates H30 for each realization. D2 and the stochastic comparator have the same source interval-calibration roles and bootstrap multiplicities. Point estimates are medians; MAE, CRPS, 90\% interval coverage and interval score are reported on available cases, with paired comparisons on common cases.

In the twenty historical source fits, we apply the same observed-window support projection to both methods after their original calibration. The seven stored quantiles do not determine their outer tails; for comparable CRPS integration we extend the tails using adjacent normal-quantile spacings and enforce $L_y$ when available. The original model weights and calibration offsets remain fixed in this source re-analysis. In the whole-archive study described below, offsets for both methods are instead re-estimated on their original independent calibration roles after projection. Stored stochastic bridge-repair tables are required to reproduce the historical source, open-source archive and group-maximum sensitivity calculations; these dependencies were checked retrospectively. Neither approach makes the assumed tail shape a verified full distribution.

The conditional-mean mechanism is useful but limited. For observed information $X$ and 30-day averaging vectors $w_j$, convexity gives
\[
 \max_j w_j^\top E(Y\mid X)\le E\{\max_j w_j^\top Y\mid X\}.
\]
It concerns the true conditional mean and conditional expectation; an approximate learner need not underestimate every individual year. Calibration, local adaptation, model bias and the chosen point summary remain empirical questions.

We evaluate a rule that reports an estimated value, an observed lower bound or an unknown state. Numerical release requires an available D2, at least four complete historical years, visible features within source-derived limits and a raw individual classifier score no greater than 0.10. The score was trained for the original unprojected D2 and is not a calibrated group-risk certificate for the projected value. A seven-day close match between visible water and air temperature is tested for the imposed exact-air-substitution scenario; that detector is not evaluated for other stresses. No observed bound is released when this quality flag is raised. The primary factorial trend comparison is separate. A second exploratory analysis passes released values and bounds into a Gaussian working likelihood, treating released numerical estimates as exact and each bound $L$ as a right-censored response with contribution $P(H30\ge L)$. Unknown cases remain in availability denominators. The resulting regression interval includes working-model variability and is not a task-calibrated missingness-only interval.

\subsection{Whole-archive degradation, common cohorts and task-level uncertainty}
The whole-archive benchmark jointly masks selected gaps in several years of a target station before fitting its daily models, climatology or residual histories. Unmasked target observations and neighboring physical stations remain available in this experiment. The benchmark contains 42 stations with at least twenty complete reference years in seven regions. One Finnish record ends in 1994, leaving 41 stations for forward evaluation. Twenty paired mask repetitions use both empirical calendars and controlled 30--90-day gaps; after seven observed burn-in years, each reference year is degraded with probability 0.35. These 3,320 archive experiments do not reuse any masked target day when fitting another year. The complete reference, including those hidden days, is used for scoring; its unmasked values can also enter local fitting. The archived single-year plug-in is retained as an optimistic-information paired comparator on identical masks.

Offline daily fits use the remaining visible history. Forward annual models and interval offsets use source observations before 2010 and exclude the evaluated test stations and their connected groups. The original partitioned role assignment can retain other training-only groups in the same named region: this occurs in the HUC17 estimator and in NVE interval calibration. Accordingly this forward experiment is held-test-group, not uniformly held-region, and is applied to 2010 onward. The strict-region sensitivity removes the additional same-region groups and is reported separately. Forward local daily fitting and dependence estimation also exclude water observations from 2010 onward. These are current-vintage retrospective forward experiments, not a historical real-time field trial. Annual regression weights remain source-only and use source-fit seed zero throughout. Source observations do not change across archive masks, so those fits are reused; every local stage that uses the degraded target archive is refitted. Twenty archive masks are not twenty new global model fits.

On each degraded archive, the retention set (90\% warm-season rule or all supported years) is crossed with the annual estimator (daily plug-in, equally calibrated stochastic regression or D2). Every slope is scored against the same complete-record slope. For the four MSE cells $A$ (rule/plug-in), $B$ (all/plug-in), $C$ (rule/alternative) and $D$ (all/alternative), the symmetric retention contribution is $\{(A-B)+(C-D)\}/2$, and the estimator contribution is $\{(A-C)+(B-D)\}/2$. They sum to $A-D$; the interaction is $(A-B)-(C-D)$. Complete paired cells, rather than different available subsets, determine each decomposition. The retention figure uses common station-repetition rows: in the forward controlled-gap setting, both retention sets are available for 807 of 820 settings, and the all-year RMSE is recomputed on those 807 rows.

We refit both methods' annual quantile offsets on their original disjoint source calibration groups after imposing the observed-window lower bound. Both distributions and their samples respect the bound when one is available. Seven stored quantiles identify only the central 90\% of the distribution. We extrapolate both outer tails using the adjacent normal-quantile spacing and project the annual samples onto the bound; the endpoint-constant alternative is retained on paired cases. The normal extrapolation is an assumption, not a validated tail law. That projection is applied to the annual distribution and then carried into the slope. It does not project a trend interval onto an annual lower bound. The main analysis uses 256 stratified samples per year. Within one estimator, independent ranks, auxiliary shared year-bootstrap regression states and three-year residual blocks keep that estimator's annual marginals fixed. The two estimators do not share one annual distribution. Shared states come from a small harmonic/air Ridge model, not the D2 tree posterior or the selected donor model. These couplings are dependence sensitivities, not posterior samples or coverage theorems.

The joint-rank numerical sensitivity fixes the annual quantiles, observed bounds and evaluation years and increases the stratified sample count through 256, 1,024 and 4,096 on the difficult forward controlled gaps. Three random seeds index the rank construction. The rank sequence is generated at 4,096 draws and truncated for smaller counts; annual marginals are regenerated separately for each sample size. We report coverage, width, interval score, the three-seed mean and the largest single-seed change. A 0.3-point rule on the three-seed mean is only a diagnostic: averaging seeds can cancel opposing changes, and meeting it does not justify ranking small construction differences. The check uses the same opened data and is not an independent validation.

Table~\ref{tab:method-versions} separates the calibration and projection order of the historical, whole-archive, task-calibration and external analyses. The task-level comparator calibrates absolute slope error divided by an independent-draw interval half-width. Every source connected group receives equal total weight in the empirical 95th percentile. Calibration uses only unused rule-calibration groups, re-predicted with the held-region model; forward errors end before 2010. We withhold the calibration when fewer than three nondegenerate source groups remain, the calibration factor is effectively zero, or the target half-width lies outside the source 1st--99th percentile range. These thresholds are empirical support checks rather than a finite-sample guarantee. Source support, availability and matched-case width, coverage and interval score are reported. All uncertainty scores target the finite complete-record slope, not climate-trend uncertainty.

The historical design metadata are not rewritten. The saved 3,320 \texttt{design.json} files state that held region and connected groups were excluded; that is correct for held test groups but not an unconditional region exclusion. The forward HUC17 estimator retained one other same-region training group, and NVE interval calibration retained two. The paper uses the held-test-group analysis as primary and the strict-region sensitivity as a separate check. A correction manifest accompanies the historical files.


\begin{table}[t]
\centering
\fontsize{8}{10}\selectfont
\renewcommand{\arraystretch}{1.08}
\newcolumntype{M}[1]{>{\raggedright\arraybackslash}p{#1}}
\caption{Calibration and projection order in the analyses used here. The rows are separate analyses; they are not twenty repetitions of one final pipeline.}
\label{tab:method-versions}
\begin{tabular}{M{0.18\textwidth}M{0.22\textwidth}M{0.22\textwidth}M{0.28\textwidth}}
\toprule
Analysis & Projection and tail treatment & Calibration and source roles & What it supports \\
\midrule
Twenty historical single-year source repetitions & The same observed-window projection is applied to both methods after their original calibration. The seven stored quantiles are extended with adjacent normal-quantile spacings. & Original weights and offsets remain fixed. Source roles are the original held-region and held-group assignments. & Annual H30 comparisons and the original H1/H2 evidence. Not a whole-archive trend refit. \\
Whole-archive factorial & The observed-bound projection is applied to each estimator's annual distribution before and after re-estimated offsets, and those annual samples then enter the slope. Each estimator has its own bounded annual marginals. Within one estimator those marginals stay fixed across the three rank constructions. Both use normal-spacing tails. A trend interval is not projected onto an annual lower bound. & Annual weights stay source-only; local daily models, climatology and residual histories are refitted after jointly masking the target archive. The main analysis excludes held test groups; the strict-region sensitivity removes additional same-region groups. & The retention/estimator decomposition and the finite-record trend comparison. The primary retention contrast uses common station-repetition rows. \\
Forward task calibration & The observed lower bound is imposed on the annual distribution. The independent-draw slope interval is then compared with the task-calibrated slope interval on the same cases. This does not project a trend interval onto an annual lower bound. & Calibration uses unused rule-calibration groups, equal group weights and pre-2010 source errors. Fewer than three nondegenerate groups, a zero factor or an out-of-range target width returns an unsupported status. & The negative result that task calibration is available for only a small selected subset and does not establish population coverage. \\
External exploratory comparisons & Source models and source calibration are fixed; local history is refitted where the adaptation requires it. & Reference quality, clock, quality codes and watershed dependence are unresolved. & Descriptive transfer evidence only. These rows cannot substitute for a strict independent reference. \\
\bottomrule
\end{tabular}
\end{table}


Two controlled-gap sensitivities assess missingness mechanisms on the first five masks at each of 41 forward-eligible stations. One hides the target's exact gaps from the donor panel. The other keeps each archive's number of missing years fixed while preferentially selecting years with high air-temperature H30. Gap locations and numbers of missing days can still change. Both tests are retrospective and do not represent independent climates or observed outage frequencies.

\section{Results}
\subsection{What the observations establish}
\input{figure_observed_records.tex}
At \RRecordStations{} of \RRecordDenominator{} eligible stations, \RRecordStationYears{} incomplete years contain an observed window above the maximum complete-year H30 (Figure~\ref{fig:observed-records}). Six stations have a witness exceeding the complete-year record by more than 0.5$^{\circ}$C. Another incomplete year could still be warmer, and measurement error or a changed sensor location can affect physical interpretation.

The practical consequence depends on the rule. Table~\ref{tab:rule-accounting} accounts for all \CCEligibleYears{} eligible station-years. Complete-year selection admits \CCCompleteAdmitted{} years and discards 561; 545 of the discarded years still contain a fully observed 30-day window and 16 are unknown. The station-specific warm-season rule admits \CCWarmAdmittedYears{} years and excludes 66, of which 50 still contain a fully observed window. Of the 35 selected witnesses, \CCWarmWitnessAdmitted{} pass the warm-season criterion and \CCAnnualWitnessAdmitted{} pass 95\% annual coverage. Among the warm-season exclusions, only \CCWarmExcludedFacts{} exceed the largest reported window among the other rule-accepted years, by \CCWarmIncrementOne{} and \CCWarmIncrementTwo{} $^{\circ}$C. The larger 0.930$^{\circ}$C difference for the Bavarian 2022 case is relative to the complete-year record, not to the warm-season accepted record; the two comparisons must not be mixed.

The rule also admits \CCWarmAdmittedNewFacts{} incomplete years whose observed windows exceed the accepted record of the other accepted years. Strict completeness discards all of them. This is the information at stake. It remains reported-window information: a rule-accepted incomplete year may have an unobserved higher H30. At an assumed absolute error of \CCMarginHalfWidth{}$^{\circ}$C on each window, both excluded differences exceed $2\epsilon$; at 0.10$^{\circ}$C, neither difference exceeds $2\epsilon$. These margins are sensitivity assumptions, not measured instrument errors.

\input{rule_accounting_table.tex}

The seasonal-window, past-only and air-regression diagnostics are reported in the Supporting Information. They show that retained heat information depends on the window definition, the historical-support requirement and regulated rivers. They do not reproduce the multiyear climate correction of \citeA{beaufort2022}, whose target and site exclusions differ, and they do not establish a universal H30 error baseline from the annual-mean analysis of \citeA{johnson2021}.

The additional observation-only USGS intake is reported in the Supporting Information. It corroborates omission at three stations but does not estimate incidence, verify a unique warmest year or establish external reconstruction accuracy.

\subsection{Annual point and distribution comparisons}
On controlled gaps overlapping at least half the training climatological peak, the median source-repetition H30 MAE is approximately 0.428$^{\circ}$C for the daily plug-in, 0.355 for functional debiasing, \WADTwoMAE{} for D2 and \WAStochasticMAE{} for calibrated stochastic regression (Figure~\ref{fig:estimators}). Median 90\% interval coverage is \WADTwoCoverage\% and \WAStochasticCoverage\%, respectively, with CRPS \WADTwoCRPS{} and \WAStochasticCRPS{}. On common cases D2 improves CRPS in \WADTwoCRPSWins{} of twenty repetitions and upper-tail CRPS in \WADTwoTailWins{}; equality of the methods is not established. Both distributions use the same support projection and tail approximation in this comparison; method-specific availability and paired counts are reported in Supporting Information Table S1. The practical evidence is a gain over the simple plug-in, not a demonstrated universal need for D2.

\input{simple_research_results.tex}
\input{simple_research_tables.tex}

\input{peak_visibility_main_table.tex}
\input{rank_convergence_table.tex}

\subsection{Source roles and missingness mechanisms}
A seed-zero audit of \RHModels{} source models found no held test station or connected group in a source role. The forward HUC17 estimator nevertheless used \RHHucCases{} cases from \RHHucGroups{} other group in the named region, and NVE interval calibration used \RHNveCases{} cases from \RHNveGroups{} other same-region groups. These are not held-target observations, but the original forward design is not uniformly held-region.

After removing the affected same-region groups and refitting, D2 trend RMSE on \SRPairN{} common forward controlled-gap archives changes from \SROldDTwo{} to \SRNewDTwo{} $^{\circ}$C per decade; stochastic RMSE changes from \SROldStoch{} to \SRNewStoch{}. A separate candidate-provider sensitivity refits all 13 source models with target-region groups excluded. On \OpenSourcePairs{} common forward controlled-gap archives in \OpenSourceGroups{} connected groups, original D2 and stochastic RMSE are \OpenSourceOldDTwo{} and \OpenSourceOldClassic{}, versus \OpenSourceNewDTwo{} and \OpenSourceNewClassic{} for the candidate models. Providers, role assignments and source sample sizes change together, so differences cannot be attributed to one factor. These analyses do not establish independent-climate transfer or resolve derivative permissions.

On \ESArchives{} forward controlled-gap archives at \ESStations{} stations in \ESGroups{} connected groups, hiding the same target gaps from the donor panel increases D2 trend RMSE from \ESBaseline{} to \ESDonor{} $^{\circ}$C per decade. The corresponding strict-original source comparison is \SDSStrictBase{} to \SDSStrictDonor{}, and the candidate-provider comparison is \SDSCandidateBase{} to \SDSCandidateDonor{} on the same \SDSCases{} common archives. In a separate heat-associated sensitivity, each archive keeps its original number of missing years: D2 RMSE is \ESHeat{}, but the number of air-hot masked years changes from \ESHotBaseline{} to \ESHotHeat{} and hidden days from \ESDaysBaseline{} to \ESDaysHeat{}. These simulations do not isolate heat association, quantify natural outage frequency or establish transfer to new climates. Supporting Information reports the pilots, both estimators, availability, intervals and group sensitivity.

In a further post-hoc sensitivity, the heat-selected years retain the baseline archive's entire set of gap offset--length templates and its hidden-day count. Among \HGArchives{} settings available for both methods and all three designs, D2 trend RMSE is \HGDBase{}, \HGDPrevious{} and \HGDMatched{} $^{\circ}$C per decade for baseline, prior heat and matched-gap heat masks, respectively; stochastic RMSE is \HGSBase{}, \HGSPrevious{} and \HGSMatched{}. Paired group-bootstrap intervals for matched minus baseline MSE cross zero for both methods. Re-scoring the matched-gap archives with the strict-original and candidate-provider source models gives D2 baseline / stress RMSE of \HGRDRBase{} / \HGRDRStress{} and \HGRDCBase{} / \HGRDCStress{} on a separate \HGRCases{}-archive common cohort. The post-hoc paired difference between candidate and original D2 stress effects is \HGRInteraction{} with group-resampled interval [\HGRInteractionLow, \HGRInteractionHigh] $({}^{\circ}$C per decade$)^2$. Its upper limit is close to zero and \HGRInfluenceCross{} of \HGRInfluenceGroups{} leave-one-group intervals cross or touch zero. Resampling was revised after inspection, with no multiplicity correction across 16 reported cells. Provider, source sample and roles change together in that comparison. These fixed source models and selected archives do not establish that heat-associated outages improve or worsen recovery; gap timing relative to each year's peak also changes.

\subsection{Release and downstream selection}
The release logic is shown in Figure~\ref{fig:release-system}; Figure~\ref{fig:output-boundaries} reports the cost of withholding numbers across the stress conditions. A score-only cutoff chosen in disjoint source rule-calibration groups approximately matches the three-output rule's source publication rate. On held-region controlled gaps, their numerical retention is \SRCGapThreeRetention\% and \SRCGapScoreRetention\%, with selected error \SRCGapThreeRisk\% and \SRCGapScoreRisk\%, respectively (Supporting Information, Table S2). Target retention differs. A separate target-distribution match yields \MRFutureCount{} numerical outputs per rule on forward controlled gaps, with failure rates \MRFutureThree\% and \MRFutureScore\% (Table S3). It uses no target outcomes to choose the cutoff, but is transductive and does not equalize bound information. In the imposed air-substitution calendar only \MRAirCount{} numbers remain; withholding nearly every value does not establish general sensor-fault detection.

Matching numerical counts separately within \AMCases{} forward controlled-gap archives yields \AMNumerical{} estimates for each rule, with \AMBadOriginal{} and \AMBadScore{} annual errors above 0.5 $^{\circ}$C. The released-values-plus-bounds working-likelihood trend RMSE is \AMOriginalRMSE{} and \AMScoreRMSE{} $^{\circ}$C per decade. A connected-group bootstrap interval for the paired squared-error difference spans zero. Selected years and bound values can differ, and this target-count match is not a prospective threshold.

A classifier trained for the support-projected D2 point also fails to show a stable advantage at matched numerical retention. On \PRCases{} forward controlled-gap masked years, source-derived old and new score thresholds retain \PROldRetention\% and \PRNewRetention\%, with selected failure \PROldFailure\% and \PRNewFailure\%. Matching \PRMatchedCount{} numerical outputs separately by target region produces \PRMatchedOldBad{} and \PRMatchedNewBad{} failures; the bootstrap interval for the failure-rate difference (new minus old) is [\PRMatchedLow, \PRMatchedHigh] percentage points. Both analyses are post-hoc and neither certifies group risk or calibrated probabilities (Supporting Information).

After inspecting reservoir inputs, we also assessed a marginal training-range warning on previously opened source rule-calibration groups. It marked 42 of 868 D2 failures among 8,010 eligible controlled-gap cases. At the same post-hoc numerical count, score-plus-warning retained 268 failures versus 211 for score-only ranking. This comparison omits quality and other range gates and chooses the matched cutoff from the evaluated score distribution; it does not establish a prospective release rule (Supporting Information).

\subsection{External reference qualification}
Strict independent-reference accuracy remains unresolved. Table~\ref{tab:external-qualification} keeps complete or bounded reference years, stations meeting each history rule, quality and clock, independence, and scored cases in separate fields. Nine Greater Wellington sites stayed at the metadata screen because temperature series were not requested, so their complete hot-year count is unknown \cite{greaterwellington2016quality}. The saved screens found zero strict complete years at 24 NEON sites and zero complete July--June thermal years at five R\'eunion sites; eight new USGS candidates contain \USGSCompleteYears{} complete station-years, but \USGSFourYearStations{} stations reach four years \cite{marzolf2025neon,hubeautemp2026,usgsstatistics2026}. That USGS count is not an hourly physical reference and does not prove independence. ECCC Red River shares a river with a source station and retains grade and time-basis questions \cite{ecccautomated2024,usgsred2025}. An Australian candidate archive had prior evaluation exposure: \AUFirstEvaluable{} of \AURefYears{} candidate station-years were previously evaluable, including all \AUFormerTestEvaluable{} designated test years \cite{bom2026}. These exclusions do not establish that the observations are erroneous.

\input{external_qualification_table.tex}

Two exploratory comparisons bound what external evidence can show. British Columbia has \BCStations{} stations and \BCYears{} station-years. Of \BCMaskTotal{} masks, \BCCases{} are controlled gaps and \BCCalendarMasks{} are calendar gaps. Plug-in, D2 and calibrated stochastic H30 mean absolute errors on the controlled gaps are \BCPluginMAE{}, \BCDTwoMAE{} and \BCStochasticMAE{} $^{\circ}$C \cite{moore2024bc,callahan2025air2stream}. Three stations share the Fraser basin; daily means may use only 20 hourly observations, their day convention is unresolved, and source quality processing used water--air comparisons \cite{wsc_temperature_faq}. Nova Scotia covers \NSExpStations{} stations, \NSExpYears{} station-years and \NSExpCases{} masks selected after inspecting availability; D2 and stochastic mean absolute errors are \NSExpDTwoMAE{} and \NSExpClassicMAE{} $^{\circ}$C \cite{novascotiawater2026}. Neither comparison is strict independent-climate validation.

The continuous USGS support at Rio Mora is the clearest warning about target definition. Three years have at least one approved unqualified continuous observation in every fixed-standard-time hour. Only \RioCompleteQuarterYear{} has every quarter-hour slot. The continuous targets for \RioBoundedYears{} assign unknown quarter-hour values to the conditional range $[-5,50]\,^{\circ}$C; they are not fully verified physical years. The \RioHourlyMasks{} hourly-supported masks give D2 MAE 0.180$^{\circ}$C against the published daily target. Using the original histories, D2 interval coverage is 53/54 against published daily H30 and 33/54 against that conditional continuous target; qualifier-aware histories yield 48/54 and 28/54. The difference can combine clock, rounding, corrections, daily preprocessing and the bound assumption. It does not identify a single source error or justify a universal physical calibration. Truckee scores \TruckeeCases{} controlled-gap masks on one station and four reported years. The reservoir comparison scores \ResScoreCases{} exploratory bounded-reference cases at \ResScoreStations{} stations and \ResScoreYears{} target station-years. Both remain descriptive checks of bounds and clock assumptions, not independent accuracy tests.

\section{Discussion}
An observed 30-day window can show that a complete-year-only reported record omits heat information without an imputation model. Most witnesses nevertheless pass the warmer-season screen, and the usable information depends on the rule and the target statistic. The seasonal and past-only sensitivities reinforce the need to specify window boundaries and historical support. They use processed current-vintage records; neither numerical agreement nor the observed bound establishes sensor accuracy.

Retaining supported years can improve the slope of the particular complete record used as reference, conditional on long-record sites, masks and available donors. The 42 stations comprise 18 connected groups, including 24 Bavarian stations. Removing Bavaria changes the fixed-prediction forward controlled-gap D2-minus-stochastic MSE difference from \SRBayernCrossOld{} under the original source model to \SRBayernCrossNew{} under the strict-source sensitivity, on \SRBayernCrossN{} common station--repetition cases. These small differences are sensitive to source and region; group resampling does not establish independent climates. A finite-record slope is not the climate-process trend.

Annual calibration, tail representation and cross-year dependence remain distinct uncertainties. On \UCCommon{} of \UCAll{} forward controlled-gap archives available to D2, which is not the two-method intersection, incomplete-year-weighted D2 annual 90\% coverage is \UCAnnual\% in \UCGroups{} connected groups. The paired comparison uses the intersection reported with Figure~\ref{fig:joint-uncertainty}. With matched ranks, constant versus normal-spacing tails give nominal 95\% trend coverage \UCTailConstant\% and \UCTailNormal\%; in a separate fixed-marginal experiment, independent, auxiliary shared-parameter and residual-block ranks give \UCIndependent\%, \UCShared\% and \UCBlocks\%, respectively (256 samples; Supporting Information, Table S4). The 4,096-sample check in Table~\ref{tab:rank-convergence} keeps coverage at 87.4--89.9\%. The two experiments use different draws, so these are not additive causal contributions. The shared states come from an auxiliary Ridge model, not the fitted D2 posterior. Within one estimator the three rank constructions keep that estimator's annual marginals fixed. Increasing the joint-rank sample count does not remove the forward undercoverage. The largest three-seed-mean coverage change is \CCRankChangeMax{} percentage points and the largest single-seed change is \CCRankSeedChangeMax{} percentage points; the mean is not a resolution guarantee, and the constructions are not ranked.

Support projection can also make annual intervals degenerate: among \SPCalendarCommon{} common forward empirical-calendar cases, D2 has \SPCalendarDTwoZero{} zero-width intervals and the stochastic method has \SPCalendarClassicZero{}; these match the reported reference but do not establish zero sensor uncertainty. In the \PVCrossHidden{} of \PVCrossAll{} original forward controlled-gap cases where the complete reference exceeds the visible bound, nominal 90\% coverage is \PVCrossDTwoCoverage\% and \PVCrossClassicCoverage\% for D2 and stochastic reconstruction. This hidden-reference stratum is unavailable at decision time. On the \OpenAnnualCases{} four-way common cases in that stratum, candidate-source coverage is \OpenAnnualDTwoNew\% and \OpenAnnualClassicNew\%; changed sources and roles preclude a single-factor attribution.

Source-only interval expansion illustrates a coverage--precision trade-off. On \SYAllSupported{} of \SYAllEligible{} common forward controlled-gap cases with both empirical corrections supported, D2 coverage changes from \SYDTwoBase\% to \SYDTwoExpanded\%, with width \SYDTwoWidthBase{} to \SYDTwoWidthExpanded{} $^{\circ}$C and interval score \SYDTwoScoreBase{} to \SYDTwoScoreExpanded{}. A more conservative source-group-maximum expansion covers \GMAnnualDNew\% of \GMAnnualCases{} common cases, but D2 width increases from \GMAnnualDWidthBase{} to \GMAnnualDWidthNew{} $^{\circ}$C and score from \GMAnnualDScoreBase{} to \GMAnnualDScoreNew{}. The two comparisons have different support sets. On \GMTrendCases{} common nontrivial archives, the same wider endpoints and independent ranks give \GMTrendDNewCov\% empirical 95\% trend coverage with tapered quantiles, while width rises from \GMTrendDBaseWidth{} to \GMTrendDTaperWidth{} $^{\circ}$C per decade and score from \GMTrendDBaseScore{} to \GMTrendDNewScore{}. Equal annual endpoints with different interior assumptions yield different trend widths. Neither source single-year masks nor the available groups establish exchangeability or a calibrated joint distribution.

The decision results similarly do not certify risk for every river. Exact-air substitution is easier to recognize than unknown field faults, while groundwater, flow, regulation and deployment can change water--air relationships. New high-quality reference measurements are needed to evaluate natural outages and transfer beyond the present archives. The three-output rule improves on the original reconstruct-all comparison only in the specified cost cells; it does not beat simple score ranking at a matched numerical count, and a nominal working interval with 99.8\% coverage has a width of 3.26$^{\circ}$C per decade. Coverage alone is not a sufficient measure of usefulness.

\section{Conclusions}
Incomplete river-temperature records can preserve observed sustained-heat information, and retaining years reduced finite-record trend RMSE relative to the evaluated warm-season screen in the tested degradation settings. The benefit of rejecting or retaining a year depends on the actual completeness rule and the intended scientific quantity. A well-calibrated classical reconstruction remains a strong annual comparator. Forward trend intervals, task-level calibration and three-output release rules require their own evidence; the present results do not certify universal numerical recovery, group safety or independent-climate transfer.

\section*{Open Research}
The source observations, processing code, case-level results and figure scripts require a permanent repository and citations to be supplied by the authors. A local review candidate includes processed daily values from 95 of the 96 previously eligible source stations, permitting an independent arithmetic audit of all 35 observation witnesses and their screening-rule decisions. The LUBW station is omitted pending permission, so the full eligible population is not reconstructed. The candidate also includes ten additional USGS daily snapshots and their observation-only replay; neither check reconstructs the main models or verifies subdaily sampling. The Supporting Information lists the access status of each external archive and the derived-file inventory.

The source includes observations from LUBW; its reuse and redistribution terms require provider-specific resolution rather than an assumed open licence. Although all whole-archive target stations have recorded reuse paths, \RightsEstimatorLUBW{} of \RightsModelCombinations{} source-model estimator roles and \RightsCalibrationLUBW{} interval-calibration roles include positive-weight LUBW cases; the access and sharing conditions of the resulting models and case-level derivatives remain unresolved. Author, affiliation, funding, conflicts of interest, permitted reviewer access and permanent deposit information remain to be supplied before submission.

\begin{figure}[p]\centering\includegraphics[width=\textwidth]{figures/simple_research/02_annual_comparators.pdf}\caption{Annual estimation under difficult controlled gaps. Both historical source-calibrated methods receive the same observed-bound projection and normal-spacing tail model. Bars show empirical repetition ranges within the same archive; these are not independent climates.}\label{fig:estimators}\end{figure}
\begin{figure}[p]\centering\includegraphics[width=\textwidth]{figures/simple_research/03_archive_factorial.pdf}\caption{Whole-archive trend errors and the symmetric retention/estimator decomposition. Left: D2 on common station-repetition rows where both retention sets are available; annotations give n and station/group counts. The forward controlled-gap contrast uses 807 of 820 settings. Right: the left bar of each pair is D2 and the right bar is the calibrated classical method. Each bar uses that method's own complete four-cell count, printed under the setting; the two counts are not one shared denominator.}\label{fig:whole-archive}\end{figure}
\begin{figure}[p]\centering\includegraphics[width=\textwidth]{figures/simple_research/04_joint_uncertainty.pdf}\caption{Left: paired annual coverage on the two-method intersection, weighted by incomplete years. Ticks are station-repetition means on the same intersection. The tick labels give the paired archive count and each method's availability out of all archives. The dashed line is the nominal 90 percent annual target. Right: nominal 95 percent trend coverage on each method's own available archives, with that archive count. The right panel is not a paired method comparison. Group-bootstrap intervals use connected groups. No gray band is shown.}\label{fig:joint-uncertainty}\end{figure}
\begin{figure}[p]\centering
\includegraphics[width=\textwidth]{figures/architecture/fig2_release_actual.pdf}
\caption{The evaluated three-output decision and exploratory trend analysis. The cap applies to an individual raw classifier score, not a certified group risk. The seven-day water--air check was assessed only for the clean and imposed exact-air-substitution cases. The selected-trend working likelihood treats released values as exact and observed bounds as censored.}\label{fig:release-system}
\end{figure}
\begin{figure}[p]\centering\includegraphics[width=\textwidth]{figures/simple_research/05_output_boundaries.pdf}\caption{Left: three outputs for empirical-calendar masks. Right: numerical-publication rates and selected errors for controlled gaps. Source diagnostics do not remove extrapolation failures. The air-agreement flag addresses the specified exact-substitution stress only; zero publication is shown as unknown rather than zero-risk success.}\label{fig:output-boundaries}\end{figure}

\bibliography{references_tr,references_review_completion}
\end{document}
